\begin{proposition}[Obstruction in the fixed log-torus category]\hypertarget{tag-030T}{}\label{push:src:B.4.24}
\label{push:B424}
\StatementRecord{push:src:B.4.24}{proved}{030T}{\SeriesRef{push:src:B.4.18}, \SeriesRef{push:src:B.4.19}, \SeriesRef{push:src:B.4.9}, \SeriesRef{push:src:B.4.30}, \SeriesRef{push:src:B.4.33}; \citet{KKN2008Root}; \citep[Definition 1.2.3, pp.~77--78]{KKN2008Root}; \citep[Definition 1.2.6, p.~79]{KKN2008Root}; \citep[Sections 1.3.1--1.3.4, pp.~81--82]{KKN2008Root}; \citet[Sections 7--9]{KKN2008PicardE}}

The local Jacobian admits a nondegenerate logarithmic
complex-torus extension over the full $(q,r)$ bidisc and hence
has proper analytic models in the sense of
\citet{KKN2008Root}. After the normalized base change
$q=Q^{12}$, the two marked axial fan presentations of
\ref{push:src:B.4.18}--\ref{push:src:B.4.19} determine a torsor under its associated
log-torus sheaf on the punctured bidisc. This torsor does
not extend across the origin as a torsor locally trivial
for the usual analytic topology. The obstruction persists
after every further finite coordinate-root base change.

On each axial chart its extending origin identifies the associated sheaf with the Jacobian log-torus sheaf $\mathcal A$. On the overlap the two identifications differ by the section $F$ in \eqref{pub:I:eq:B-4-24-3}. An extension here is a locally trivial $\mathcal A$-torsor carrying these prescribed axial identifications.
\end{proposition}



\begin{proof}
\begin{enumerate}[label=\arabic*.,leftmargin=*]
\item \emph{The Jacobian pairing.}
Retain the period columns and logarithms of \ref{push:src:B.4.18}:
\[
 \gamma=(6\mu,C-s-12v)^t,\quad
 u=(\tau,\mu)^t,\quad w=(1,0)^t,\quad\delta=(0,1)^t,
 \qquad
 s=\frac{\log q}{2\pi i},\quad
 v=\frac{\log r}{2\pi i}.
\]
Here $\tau-s,\mu,C$ are holomorphic at $q=0$.
Use multiplicative characters
$Z=\exp(2\pi i\zeta_1)$ and
$W=\exp(-2\pi i\zeta_2)$, with character lattice
$X=\Z e_1\oplus\Z e_2$ and period lattice
$Y=\Z u\oplus\Z\gamma$. Put
\[
 a_0=e^{2\pi i(\tau-s)},\qquad
 a_1=e^{2\pi i\mu},\qquad a_2=e^{-2\pi iC}.
\]
These are holomorphic units on the full $q$ disc.
The multiplicative pairing is
\begin{equation}
 \begin{array}{c|cc}
       &u&\gamma\\\hline
 e_1&q a_0&a_1^6\\
 e_2&a_1^{-1}&q r^{12}a_2.
 \end{array}
\label{pub:I:eq:B-4-24-1}
\end{equation}
Give the base the fine saturated log structure from
$qr=0$. In its sharp monoid $\N^2$, written additively,
the pairing modulo units is $\operatorname{diag}(q,q+12r)$.
We check the exact admissibility definition
\citep[Definition 1.2.3, pp.~77--78]{KKN2008Root}.
On the full face a nonnegative functional
$N(q)=h,N(r)=k$ gives
$\operatorname{diag}(h,h+12k)$. If $h>0$ it is
nondegenerate. If $h=0<k$, its kernel on each lattice
is the first basis line, exactly the face lattice
for the face generated by $q$; the quotient pairing
is multiplication by $12k$. For the $q$ face the
nonzero face lattices are these first basis lines
and the nontrivial quotient pairing is multiplication
by $h$. The $r$ face and the zero face have zero
face lattices. Zero functionals give zero-dimensional
quotients. These are all the face tests.

\item The residual unit condition is also satisfied
\citep[Definition 1.2.6, p.~79]{KKN2008Root}.
At $q=0$ the residual lattices vanish. At $r=0$,
$q\ne0$, their pairing is $\log|q a_0|\ne0$ after
shrinking the disc. At ordinary points the real
logarithmic matrix has determinant
\[
 \log|q a_0|\log|q r^{12}a_2|
                 +6\bigl(\log|a_1|\bigr)^2>0,
\]
after shrinking so both of the first two logarithms
are negative. Thus \eqref{pub:I:eq:B-4-24-1} defines a log complex
torus $\mathcal A$ by Definition 1.3.4, pp.~81--82.
Theorems 5.2.1 and 5.3.2, pp.~171--172 and 180--181,
respectively supply a complete invariant fan with
finitely many $Y$-orbits and its proper analytic
model.

\item \emph{The two extending origins.}
Set
\[
 B=\Delta_Q\times\Delta_r,\quad q=Q^{12},\qquad
 B_q=\Delta_Q^*\times\Delta_r,\quad
 B_r=\Delta_Q\times\Delta_r^*,\quad U=B_q\cap B_r.
\]
On $B_r$, \ref{push:src:B.4.9} supplies the affine section
\begin{equation}
 \sigma=
 \left(6\mu v,\,-6v^2+v(C-12\kappa_0)\right)^t.
\label{pub:I:eq:B-4-24-2}
\end{equation}
For every local branch of $\log r$, both of its
multiplicative coordinates are holomorphic units
in the full $q$ disc. Its exact $r$-deck difference
is the bounded cusp closing translation
of \ref{push:src:B.4.9}. Hence it descends as a global section
of the affine family, including $q=0$.
Translation by $-\sigma$ identifies the prescribed
cusp fan presentation with its group presentation.
The normalized toric base change $q=Q^{12}$ preserves
this identification and its extending origin.

\item On $B_q$, the rotating model of \ref{push:src:B.4.18} is the
free cyclic quotient of its root fan model. Selecting
$Q$ recovers that model; more generally the normalized
base change recovers it also at the ramification
divisor. The twelfth fractional deck is
$\gamma-6\delta$, a period, so the
multiplicative period action is exactly the second
column of \eqref{pub:I:eq:B-4-24-1}. The point $(Z,W)=(1,1)$
extends in the root fan at $r=0$ and gives the
group origin. Thus both axial presentations have
extending origins in their associated
$\mathcal A$-sheaves. The identity on the character and period lattices identifies the associated quotient sheaves under a fan refinement: both are defined by the same multiplicative pairing \eqref{pub:I:eq:B-4-24-1}.

\item In these two origins the exact comparison
$\zeta\mapsto\zeta-\sigma+F$ of \ref{push:src:B.4.18} becomes
translation by
\begin{equation}
 F(s,v)=\left(
 \frac{\mu(s+12v)}2,
 \frac{C(s+12v)}{12}-\frac{(s+12v)^2}{24}+6v
 \right)^t.
\label{pub:I:eq:B-4-24-3}
\end{equation}
This defines the associated $\mathcal A$-torsor
by gluing trivial torsors on $B_q$ and $B_r$.
The differences are exactly
\begin{equation}
 F(s+12,v)-F(s,v)=\gamma-6\delta,
 \qquad F(s,v+1)-F(s,v)=\gamma.
\label{pub:I:eq:B-4-24-4}
\end{equation}
They are periods, so $F$ is an $\mathcal A$-section
on $U$.

\item \emph{The ghost-valuation obstruction.}
By the quotient-sheaf definition
\citep[Sections 1.3.1--1.3.4, pp.~81--82]{KKN2008Root},
on the underlying analytic site there is an exact sequence
\begin{equation}
 0\longrightarrow Y\longrightarrow
 \mathcal T_{\log}^{(Y)}\longrightarrow
 \mathcal A\longrightarrow0.
\label{pub:I:eq:B-4-24-5}
\end{equation}
The middle term is the sheaf of bounded multiplicative
log characters. On $U$ it is the ordinary
two-dimensional multiplicative torus sheaf.
In particular \eqref{pub:I:eq:B-4-24-4} says that the connecting
class of $F$ in $H^1(U,Y)$ has $r$-loop component
$\gamma$. The terms in $\Z w+\Z\delta$ disappear
on exponentiation.

\item Let $i:D_Q=\{Q=0\}\hookrightarrow B_r$.
Taking the $Q$-order of each character gives
\begin{equation}
 \operatorname{ord}_Q:
 \mathcal T_{\log}^{(Y)}\longrightarrow i_*\Z^2,
 \qquad
 u\longmapsto(12,0),\quad
 \gamma\longmapsto(0,12)
\label{pub:I:eq:B-4-24-6}
\end{equation}
on its period subsheaf. If $h$ is any
$\mathcal A$-section on $B_r$, local lifts in
the middle sheaf of \eqref{pub:I:eq:B-4-24-5} represent its
connecting class by a $Y$-cocycle. Applying
\eqref{pub:I:eq:B-4-24-6} makes that cocycle a coboundary,
the difference of the local ghost vectors.
But restriction identifies $H^1(B_r,Y)$ with
$H^1(D_Q,Y)$, since $B_r$ retracts onto $D_Q$;
the induced map to $H^1(D_Q,\Z^2)$ is the
injective diagonal map $(12,12)$. Therefore
the connecting class of $h$ is zero. In
particular its $r$-loop component is zero.

\item For a section on $B_q$ that component is zero
as well, because the $r$ loop contracts in $B_q$.
Consequently $F$ cannot be the difference of
sections from $B_q$ and $B_r$. The glued torsor
is nontrivial on $B\setminus\{0\}$. An extension
locally trivial near the origin would have a
section there; after shrinking the bidisc its
two coordinate expressions would give precisely
the forbidden splitting. This proves nonextension.

\item For $q=Q^{12m}$ and $r=R^n$, with $m,n>0$,
the $R$-loop difference is
\[
 F(s,v+n)-F(s,v)
    =n\gamma-6n(n-1)\delta.
\]
Its $Y$-component is $n\gamma\ne0$,
whereas the $Q$-ghost map on $Y$ is
$\operatorname{diag}(12m,12m)$ and remains
injective. The same proof applies after every
such cover. Finite Kummer covers of the standard
coordinate log structure are locally dominated
by coordinate-root covers, so allowing them
does not remove this obstruction either.

\item \emph{The ordinary gauge.}
The exact identity in Proposition~\ref{push:src:B.4.19} is
\begin{equation}
 F-\sigma-kv\delta=G(s)+6v\delta,
 \quad k=-s+12\kappa_0,
 \quad G(s)=\left(\frac{\mu s}2,
                  \frac{Cs}{12}-\frac{s^2}{24}\right)^t.
\label{pub:I:eq:B-4-24-7}
\end{equation}
On $q\ne0$, this is the raw overlap in the ordinary gauge. Its second multiplicative character is
\[
 e^{-2\pi i kv}=Q^{12v}e^{-24\pi i\kappa_0v},\qquad q=Q^{12}.
\]
The second factor is a holomorphic unit on a fixed branch of $\log r$. A bounded logarithmic character at $Q=0$ has integral $Q$-order, while $12v$ varies with $r$. Multiplication by a period changes that order by an integer. Thus this gauge cannot give an extending log-torus section along the full $Q$-axis. Changes of the two extending axial origins instead lie in the image of
\[
 H^0(B_q,\mathcal A)\oplus H^0(B_r,\mathcal A)
 \longrightarrow H^0(U,\mathcal A),\qquad(h_q,h_r)\longmapsto h_r-h_q.
\]
The ghost calculation shows that $F$ lies outside this image.

\end{enumerate}
\end{proof}